\section{Estado del Arte}\label{sec:eda}

La revisión de la literatura en torno al Problema del Clique Máximo (MCP) revela un esfuerzo de más de seis décadas por encontrar soluciones algorítmicas eficientes a este desafío NP-duro \cite{bomze1999, pardalos1994maximum}. A lo largo de los años, los enfoques han evolucionado desde la combinatoria discreta pura hacia técnicas de optimización continua y heurísticas avanzadas. A continuación, se detallan los avances predominantes:

\begin{enumerate}
    \item \textbf{Complejidad y Modelos Matemáticos:} El MCP es computacionalmente equivalente al problema del Conjunto Independiente Máximo (Maximum Independent Set) y al problema de Cobertura de Vértices Mínima (Minimum Vertex Cover) \cite{pardalos1994maximum}. El problema se modela comúnmente mediante programación entera (como un problema binario discreto). Una formulación alternativa de naturaleza continua, basada en el teorema de Motzkin-Straus, se detalla en la Sección~\ref{subsec:motzkin}. Desde la perspectiva de la complejidad, no solo es NP-duro encontrar la solución exacta, sino que se ha demostrado que es imposible aproximar el tamaño del clique máximo dentro de un factor constante en tiempo polinomial, a menos que P = NP \cite{pardalos1994maximum}.
    
    \item \textbf{Métodos de Solución Exactos:} En el ámbito de la optimización exacta, el enfoque principal ha sido el método de Ramificación y Acotamiento (\textit{Branch \& Bound}). Históricamente, el algoritmo de Carraghan y Pardalos (1990) sentó las bases algorítmicas para los solucionadores modernos, utilizando enumeración parcial y acotamiento agresivo \cite{szabo2018_different}. Solucionadores actuales de vanguardia como \textit{cliquer} (Östergård), MCR (Tomita) y BBMC (San Segundo) son descendientes y versiones refinadas de este enfoque clásico \cite{szabo2018_different}. Recientemente, se ha demostrado que reformular el problema de optimización en una serie de problemas de decisión parametrizados (búsqueda de un clique de tamaño específico $k$) simplifica la estructura del problema y permite reducir drásticamente el espacio de búsqueda \cite{szabo2018_different}.
    
    \item \textbf{Enfoques Heurísticos:} Dada la dificultad de los métodos exactos en grafos muy grandes, la literatura ha propuesto múltiples heurísticas para aproximar soluciones. Las más comunes son las heurísticas codiciosas secuenciales (\textit{sequential greedy heuristics}), que construyen cliques añadiendo el mejor vértice disponible ("Best in") o eliminando el peor ("Worst out") \cite{pardalos1994maximum}. Para mejorar estas aproximaciones, se aplican algoritmos de búsqueda local (como intercambios-$k$), búsquedas aleatorias (\textit{Randomized algorithms}), búsqueda Tabú y redes neuronales \cite{pardalos1994maximum}. Enfoques más recientes aprovechan el cálculo de Conjuntos Independientes Mínimos (Minimal Independent Sets) en cada paso recursivo, descartando tempranamente nodos que no pueden mejorar la mejor solución actual (poda), reduciendo eficientemente el árbol de exploración \cite{singh2014_simple}.

    \item \textbf{Casos Polinomiales y Grafos Especiales:} Aunque el problema es NP-duro en grafos arbitrarios, el estado del arte reconoce que puede ser resuelto en tiempo polinomial en clases de grafos con topologías especiales. Destacan de manera fundamental los grafos perfectos (\textit{perfect graphs}), grafos bipartitos, grafos de intervalos y grafos triangulados \cite{pardalos1994maximum}. En estas estructuras, la relajación lineal del problema entero entrega soluciones óptimas con valores enteros de manera natural \cite{pardalos1994maximum}.
\end{enumerate}

\subsection{Formulación cuadrática continua: Teorema de Motzkin-Straus}\label{subsec:motzkin}

Una formulación alternativa del MCP, ampliamente referenciada en la literatura, proviene del teorema de Motzkin-Straus \cite{motzkin1965}, que establece una equivalencia exacta entre el problema combinatorio discreto y un problema de programación cuadrática continua. Sea $A$ la matriz de adyacencia del grafo $G$ y $\Delta$ el simplex estándar $\Delta = \{x \in \mathbb{R}^n : \sum_{i=1}^{n} x_i = 1,\ x_i \geq 0\}$. El teorema establece que:

\begin{equation}\label{eq:motzkin}
    \frac{1}{2}\left(1 - \frac{1}{\omega(G)}\right) = \max_{x \in \Delta}\ \frac{1}{2}\, x^T A x,
\end{equation}

\noindent donde $\omega(G)$ denota el tamaño del clique máximo de $G$. Esta reformulación transforma el problema discreto en uno de optimización continua no lineal, lo cual permite emplear técnicas de programación cuadrática para su resolución \cite{pardalos1994maximum, bomze1999}. Aunque esta formulación no se implementa en el presente trabajo, su existencia evidencia que el MCP admite paradigmas de modelado más allá de la programación lineal entera, lo cual ha motivado numerosas extensiones, como relajaciones semidefinidas y variantes regularizadas del problema cuadrático \cite{bomze1999}.

\subsection{Contextos de aplicación}\label{subsec:aplicaciones}

A continuación se desarrollan dos contextos en los que el MCP aparece como modelo central, justificando su interés práctico más allá del valor teórico.

\subsubsection{Teoría de códigos y diseño de códigos correctores de errores}

El MCP surge de manera natural en la construcción de códigos binarios de longitud $n$ y distancia mínima de Hamming $d$. Se define el grafo $G_{n,d}=(V,E)$ con $V=\{0,1\}^n$ y $\{u,v\}\in E$ si y solo si $d_H(u,v)\geq d$. Un clique en $G_{n,d}$ corresponde a un conjunto de palabras código mutuamente separadas a distancia $\geq d$; el clique máximo determina el tamaño $A(n,d)$ del código óptimo, cantidad de interés fundamental en teoría de la información \cite{bomze1999}.

Este vínculo es particularmente relevante para la familia de instancias \textit{Keller} utilizada en este trabajo. Los grafos $G_n$ de Keller modelan la conjetura sobre teselaciones del espacio $\mathbb{R}^n$ por hipercubos unitarios trasladados: la existencia de un clique de tamaño $2^n$ en $G_n$ implica un contraejemplo a la conjetura para dimensión $n$ \cite{lagarias1992keller}. Resolver el MCP en estas instancias ha sido decisivo para refutar la conjetura en dimensiones $n \geq 8$ y para estudiar empaquetamientos densos relacionados con códigos.

\subsubsection{Bioinformática: alineamiento de estructuras moleculares}

En biología computacional, el MCP se utiliza para encontrar la \emph{subestructura común máxima} (Maximum Common Substructure, MCS) entre dos moléculas o proteínas. Dadas las estructuras $\mathcal{A}$ y $\mathcal{B}$, se construye un \emph{grafo de productos} $G_P$ cuyos vértices son pares $(a,b)$ con $a \in \mathcal{A}$, $b \in \mathcal{B}$ compatibles (mismo tipo de átomo o residuo), y dos vértices $(a_1,b_1)$ y $(a_2,b_2)$ se conectan si las distancias internas $\|a_1-a_2\|$ y $\|b_1-b_2\|$ coinciden dentro de una tolerancia. Un clique máximo en $G_P$ representa el subconjunto máximo de pares mutuamente consistentes, es decir, la mejor superposición posible \cite{bomze1999, gardiner1997clique}.

Esta formulación es la base de algoritmos de \textit{docking} molecular y descubrimiento de fármacos, donde identificar farmacóforos comunes entre ligandos activos guía el diseño de nuevas moléculas terapéuticas \cite{gardiner1997clique}.